\frontchapter{Abstract}
People with subjective cognitive decline, mild cognitive impairment, or
dementia often lose track of where they placed objects and whether they
completed daily tasks, which erodes independence and burdens caregivers. Among
the systems reviewed in this report, none combines a persistent, spatially
grounded memory of first-person experience with progressive temporal
compression and evidence-based natural-language querying for both the person
and a delegated caregiver. This report presents \textbf{Aegis}, an egocentric
spatio-temporal vision system that turns a wearable first-person video stream
into a multi-resolution memory of objects, layout, and activities: recent
history is kept in fine detail, older history is condensed into summaries.
Users and caregivers ask spatial and temporal questions by speech or text and
receive answers grounded in stored evidence, with an explicit ``not sure''
response when evidence is weak.

The project follows an iterative, vertical-slice approach. Iteration~1 (to the
mid-evaluation) delivered the batch reconstruction, perception and scene-graph
slice. Iteration~2 turned it into one end-to-end pipeline: live phone capture
with criterion-based keyframes, gap and pause records and crash-safe resume; a
gravity-aligned spatial model with nine object--anchor relations; debounced
object state runs; a persistent per-home SQLite memory with cross-session
matching of rooms, anchors and objects; verified three-day-then-summary
compaction; a rule-based, closed-vocabulary query engine served from Modal over
HTTPS; and patient and caregiver Flutter applications with on-device speech
input. The automated suite has 451 passing tests (79\% line coverage).
Component measurements show an open-vocabulary detection F1 of 0.721 on a
ten-frame, five-environment hand-labelled benchmark, and a 5.7--7.2$\times$
smaller memory with 100\% semantic retention on two short clips. A seven-day
public egocentric study (EgoLife) found that individual objects remained in
memory for one to four days, but class-level queries did not yet benefit from
the longer history. Temporal questions, calibrated confidence, caregiver
authentication, latency and usability remain unmeasured or planned for FYP-2.

\medskip\noindent\textbf{Keywords:} egocentric vision, spatio-temporal memory,
scene graph, natural-language querying, assistive technology, cognitive
impairment, memory abstraction, caregiver access.
